\chapter{P2：从全量数据找到适合解决的业务范围}

\section{承上启下}

P1 已经确认了一批可以重复使用的数据。现在的新问题不是“还有多少文件”，而是“先服务哪一类商品，运营问题才最清楚”。P2 因此要从全量数据中选出一个可以学习、可以检索、可以做运营演示的范围。

\section{业务问题、输入与边界}

业务问题是：哪些候选类目同时有足够商品、足够评论和足够的字段覆盖，能够支持单品诊断、竞品比较和问题分流？输入只能来自 P1 登记过的评论和商品文件。

这里的价值要先说清楚。\textbf{业务价值}是避免一开始就凭兴趣选类目，减少后续返工；\textbf{技术价值}是学会在大文件上做可重复的质量检查。P2 不证明任何商品真的更好，也不估计销量、缺陷率或商业收益。

\section{开发过程}

\subsection{1. 先确认没有偷偷换数据}

脚本先检查 P1 的数据指纹和登记清单。可以把它理解成进仓前核对“这是不是之前登记的那批货”。如果文件内容已经变化，流程停止，而不是看到同名文件就继续。

\subsection{2. 一行一行检查大文件}

评论文件很大，不能为了统计就把全部内容一次塞进内存。程序逐行读取，记录格式是否正确、字段是否存在、评分和时间是否合理、评论能否找到对应商品。这种边读边统计的方法叫\textbf{流式审计}，先理解为“大文件的流水点数”。

\subsection{3. 只用明确的商品键关联}

评论和商品资料通过冻结的父商品键连接，不通过相似标题猜测。结果显示商品资料 94,327 行、评论 2,128,605 行，其中 2,100,939 条评论被精确关联，27,666 条未关联。这里的“关联”只说明源字段一致，不是对现实商品身份的独立认证。

\subsection{4. 让算法给出建议，再由业务目标做选择}

先设置最少商品数、最少评论数和覆盖门槛，再按固定权重给候选范围排序。算法建议 Water Filters；Ice Makers 也满足门槛，并且更容易围绕制冰质量、噪声、可靠性、清洁、速度、安装、漏水和便携性组织运营问题，所以项目选择 Ice Makers。

这个选择记录为一个\textbf{业务覆盖决策}（override）：不是把算法改到“刚好选中 Ice Makers”，而是在合格候选中公开说明为什么选择它。最终范围包含 3,278 个商品和 136,604 条精确关联评论，文本覆盖为 0.998675。

\section{本阶段学到的知识}

\begin{itemize}
  \item \textbf{流式审计：}像流水线点数，适合大文件，且不制造完整解压副本。
  \item \textbf{精确关联：}像用订单号找订单，不用商品标题相似度猜测。
  \item \textbf{候选评分：}把“数据适不适合做项目”变成可以比较的规则。
  \item \textbf{业务覆盖决策：}算法建议和产品选择可以不同，但理由必须留下。
\end{itemize}

\begin{beginnerbox}
项目例子：Water Filters 分数更高，是数据规则给出的建议；Ice Makers 被选择，是带理由的产品范围决定。两者都可以是真实记录。真正的问题是事后偷偷改权重，却不告诉读者。
\end{beginnerbox}

\textbf{小练习：}如果某个类目故事很好讲，但评论覆盖没有达到最低门槛，能不能直接选？不能。业务判断可以在合格候选中做，不能绕过数据检查。

\section{问题、修复与验证}

最容易出现的错误，是把“最高分”误解成“唯一正确答案”。修复是同时保留算法推荐、合格门槛、最终选择和选择理由。验证包括重新核对数据指纹、格式和关联统计、覆盖门槛、候选排序及范围登记。

\section{阶段产出与验收边界}

产出包括\filelink{artifacts/audits/phase-2-audit.json}、\filelink{artifacts/manifests/phase-2-scope.json}和\filelink{docs/reports/phase-2-audit.md}。P2 在其范围内为\statuspassed：说明选定的 Ice Makers 范围来自合格数据且过程可追踪；不说明它的商业价值已经被证明，也不说明评论代表全部消费者。

\section{怎样进入下一阶段}

P3 要在这批 Ice Makers 数据上规定“什么可以算作证据”：例如噪声、漏水怎样分类，支持和反对怎样区分，结果怎样引用。P3 必须继续使用本阶段的范围登记，不能看到后面的结果后重新选类。

\section{面试怎么讲}

\begin{interviewbox}
“我先在固定输入上逐行审计和精确关联，再按门槛和固定权重比较候选。算法推荐 Water Filters，但我选择了同样合格、且更适合三个运营工作流的 Ice Makers，并把这个业务决策和理由单独记录，没有反向调整参数迎合结论。”
\end{interviewbox}

\textbf{掌握标准：}能分开说清楚“数据观察”“算法建议”和“业务选择”，并说明它们各自的证据来源。
